% frontmatter/firstwords/verano.tex -- instrucciones para el adulto del
% cuaderno de verano de «First Words» (ediciones/firstwords-verano.tex,
% ver notes/07-ediciones.md), en inglés como todo el cuaderno: van en vez
% de frontmatter/firstwords/como-usar.tex, que habla del año entero.
% Como en verano ya no se presentan tricky words, aquí están todas las que
% el libro entero ha presentado antes (\trickyLista; la lista la escribe
% tools/gen_palabras.py). Detrás, la tabla de sonidos (la del libro
% entero) y, en vez del mapa del curso, el mapa del verano (\mapaEdicion,
% preamble.tex).
\section*{How to use this summer book}

This book is the summer of \textit{First Words}: its last 13 weeks, as
a book of its own. It is the summer when a child goes from reading words
to reading sentences: every day, a short sentence of two to four words
(\textit{Toby sleeps.}, \textit{Hens cluck.}). It is for a child who can
already read English words sound by sound -- with the sounds on the next
page, and the tricky words below -- and is starting to put them
together. After it comes \textit{Read and Draw}.

There are 65 pages, one for every weekday; five or ten minutes a day are
enough. Every page has the same parts:

\begin{itemize}[leftmargin=6mm]
\item \textbf{Date}, \textbf{Read} and \textbf{Notes}, for the grown-up:
  tick every day how the reading went. At the end of the summer, it shows
  how far the child has come.
\item \textbf{The sentence of the day}, in the green box, big: read it
  pointing to each word. The words have no sound buttons any more; if
  one is hard, sound it out together.
\item \textbf{An activity}, which always ends with a picture to colour.
  The instructions (small and grey) are \textbf{for the grown-up to read
  aloud}: the child only reads what is printed big.
\end{itemize}

\subsection*{The activities}

{\small
\begin{itemize}[leftmargin=6mm, itemsep=1pt, topsep=2pt]
\item \textbf{\lblColorea}: a big picture of what the sentence says:
  the simplest way to check that it has been understood, not only
  sounded out.
\item \textbf{\lblCompleta}: colour the picture and draw what is missing.
\item \textbf{\lblTraza}: trace a letter, then write it.
  \textbf{\lblEscribe}: the same, with a whole word.
\item \textbf{\lblEncuentra}: circle a word among others that look very
  like it (\textit{sleeps, sheep, slips}\ldots): every letter counts.
\item \textbf{\lblAdivina}: the grown-up reads a riddle; the child
  colours the picture that answers it, and writes it.
\item \textbf{\lblUne}: join each word to its pair: opposites,
  rhymes\ldots
\item \textbf{\lblSiNo}: decide whether very short sentences make sense.
\item \textbf{\lblRelee} and \textbf{\lblRepasa} (Fridays): the sentences
  of the week, to read again and tick the ones that come on their own.
\end{itemize}}

\subsection*{Tricky words}

{\small These words don't sound the way they are written, so they are
read whole, like the child's own name. The summer sentences use them; if
one is new, read it together before the page.\par}
\vspace{2mm}
{\raggedright\input{content/firstwords/generated-tricky\sufijoEdicion}\par}

\subsection*{How to help}

{\small
\begin{itemize}[leftmargin=6mm, itemsep=1pt, topsep=2pt]
\item Stuck on a word? Say its sounds together, then let the child say
  the word. If it doesn't come, say it yourself and let them repeat it
  -- better that than a fight.
\item Say the sounds, not the names of the letters (\textit{sss}, not
  \textit{ess}), and keep them short.
\item If it is hard, read the page of the day before again. And the
  activity can wait: the reading is what matters.
\end{itemize}}
\newpage
